\subsection{Normalization and the spaces of prime distributions}
\label{v4-arith:pm:chapter}

The exponential coordinate $m=e^t$ changes the growth of a counting
measure. It also fixes the weight appearing in a Fourier formula.
We use additive Dirac masses on $\mathbb R$, the test space
$\mathscr A=C_c^\infty(\mathbb R)$ with its usual inductive-limit
topology, and the Schwartz space $\mathscr S(\mathbb R)$.
Distributions are complex-linear in their test functions.
Write $\gamma$ for Euler's constant and
\[
 \xi(s)=\xi_{\mathbb R}(s)
   =\tfrac12s(s-1)\pi^{-s/2}\Gamma(s/2)\zeta(s).
\]

\begin{lemma}[growth and zero counting from the theta integral]
\label{v4-arith:pm:zero-count}
For $R\geq2$ one has
$\log\max_{|s|\leq R}|\xi(s)|\leq CR\log(R+2)$.
Consequently the number of its zeros in $|s|\leq R$, with
multiplicity, is $O(R\log(R+2))$, and
$\sum_\rho|\rho|^{-2}<\infty$.
There are infinitely many zeros. All lie in $0\leq\Re s\leq1$.
\end{lemma}
\begin{proof}
In the theta integral \eqref{v4-arith:mo:xi-integral}, the theta
tail on $x\geq1$ is bounded by $C e^{-\pi x}$. For $|s|\leq R$
the powers of $x$ in that integral are bounded by $2x^{(R+1)/2}$.
The estimate
$\sup_{x>0}x^A e^{-\pi x/2}=(2A/(\pi e))^A$
bounds the integral by $\exp(CR\log(R+2))$; the quadratic
prefactor does not change this bound. Jensen's formula at radius
$2R$, using $\xi(0)=1/2$, bounds each zero in radius $R$ by a
contribution at least $\log2$. This proves the zero count.
Dyadic summation then proves the inverse-square assertion.

The growth bound gives entire order at most one. If there were
only finitely many zeros, Hadamard's factorization theorem would
give $\xi(s)=e^{a+bs}P(s)$ with a polynomial $P$.
For real $s\geq2$, however, $\zeta(s)\geq1$ and integration of
the gamma integral just over $[s/2,s/2+1]$ gives
$\log\xi(s)\geq(s/2)\log s-Cs$. This contradicts that finite
factorization. Finally, the Euler product excludes zeros for
$\Re s>1$, and the functional equation excludes $\Re s<0$.
\end{proof}

\begin{proposition}[the constant in the genus-one product]
\label{v4-arith:pm:constant}
For this normalization, $\xi(0)=1/2$, and the constant multiplying
$s$ in the canonical genus-one product is
\[
 B=\frac{\xi'(0)}{\xi(0)}
   =-1-\frac\gamma2+\frac12\log(4\pi).
\]
Thus, with zeros counted with multiplicity,
\[
 \xi(s)=\tfrac12e^{Bs}\prod_\rho(1-s/\rho)e^{s/\rho},\qquad
 \frac{\xi'(s)}{\xi(s)}
   =B+\sum_\rho\left(\frac1{s-\rho}+\frac1\rho\right).
\]
The factors and the paired summands converge locally uniformly away
from their indicated zeros or poles. The individual sum
$\sum_\rho1/\rho$ is not separated from its partner without a
specified summation convention.
\end{proposition}
\begin{proof}
The identity $\Gamma(1+s/2)=(s/2)\Gamma(s/2)$ removes the apparent
singularity at zero:
\[
 \xi(s)=(s-1)\Gamma(1+s/2)\pi^{-s/2}\zeta(s).
\]
The values $\zeta(0)=-1/2$, $\zeta'(0)=-\tfrac12\log(2\pi)$,
and $\Gamma'(1)=-\gamma$ give
\[
 \frac{\xi'(0)}{\xi(0)}
 =-1-\frac\gamma2-\frac12\log\pi+\log(2\pi).
\]
For example, the two zeta values follow directly by expanding
$\zeta(s)=2^s\pi^{s-1}\sin(\pi s/2)\Gamma(1-s)\zeta(1-s)$:
the constant terms cancel in
$\Gamma(1-s)\zeta(1-s)=-s^{-1}+O(s)$.
The remaining factor is
$\tfrac s2(1+s\log(2\pi)+O(s^2))$.
The exact special values are also recorded in
\href{https://dlmf.nist.gov/25.6.E1}{DLMF 25.6.1},
\href{https://dlmf.nist.gov/25.6.E11}{25.6.11}, and
\href{https://dlmf.nist.gov/5.4.E11}{5.4.11}.

Lemma~\ref{v4-arith:pm:zero-count} supplies the growth and zero bound
for the order-one factorization. On a compact set with
$|s/\rho|\leq1/2$, the logarithm of a factor and its logarithmic
derivative are bounded by a constant times $|\rho|^{-2}$.
Hence the product and paired derivative converge normally there.
Every canonical factor has derivative zero at $s=0$, so its
exponential constant is precisely the logarithmic derivative just
computed. This also specifies the normalization of the product.
\end{proof}

\begin{lemma}[an elementary lower bound for logarithmic prime mass]
\label{v4-arith:pm:growth}
Put $\psi(x)=\sum_{m\leq x}\Lambda(m)$. For every integer $n\geq1$,
\[
 2n\log2-\log(2n+1)\leq\psi(2n),
 \qquad \psi(x)\leq x\log x\quad(x\geq2).
\]
Consequently there are positive constants $c,x_0$ with
$\psi(x)\geq cx$ for $x\geq x_0$.
\end{lemma}
\begin{proof}
The largest of the $2n+1$ binomial coefficients in $(1+1)^{2n}$
is $\binom{2n}{n}$, so it is at least $4^n/(2n+1)$.
For every prime $p$, its valuation is
\[
 v_p\binom{2n}{n}
 =\sum_{k\geq1}\left(\left\lfloor\frac{2n}{p^k}\right\rfloor
                 -2\left\lfloor\frac n{p^k}\right\rfloor\right).
\]
Each summand is zero or one and vanishes when $p^k>2n$.
Taking logarithms therefore gives
$\log\binom{2n}{n}\leq\psi(2n)$.
The upper bound uses $0\leq\Lambda(m)\leq\log m$.
For arbitrary sufficiently large $x$, take $n=\lfloor x/2\rfloor$
and use monotonicity. No prime-number asymptotic is needed.
\end{proof}

\begin{theorem}[the exact exponential weight for temperedness]
\label{v4-arith:pm:weights}
For real $\sigma$, define the positive Radon measure
\[
 \mu_\sigma=\sum_{m\geq2}\Lambda(m)m^{-\sigma}\delta_{\log m}.
\]
Every $\mu_\sigma$ is a distribution on $\mathscr A$.
It is a tempered distribution exactly when $\sigma\geq1$.
It is finite exactly when $\sigma>1$; $\mu_1$ has infinite total mass.
In particular, $\mu_0=\mu_{\mathrm{prime}}$ is unchanged as a
locally finite measure, and $\mu_{1/2}$ is not tempered.
\end{theorem}
\begin{proof}
Only finitely many atoms meet any compact interval. Their total
variation bounds evaluation by the supremum norm on that interval,
which proves continuity on $\mathscr A$.
For $\sigma\geq1$,
\[
 |\mu_\sigma(\varphi)|
 \leq\sup_t(1+|t|)^3|\varphi(t)|
       \sum_{m\geq2}\frac{\log m}{m(1+\log m)^3}.
\]
The last series converges by the integral test. If $\sigma>1$,
$\sum(\log m)m^{-\sigma}<\infty$ also gives finiteness.
Partial summation gives
\[
 \sum_{m\leq X}\frac{\Lambda(m)}m
 =\frac{\psi(X)}X+\int_1^X\frac{\psi(x)}{x^2}\,dx,
\]
which diverges by Lemma~\ref{v4-arith:pm:growth}.

If a positive measure supported in $[0,\infty)$ is tempered, choose
a nonnegative smooth compact cutoff equal to one on $[0,1]$ and
rescale it by $t\mapsto t/T$. The finite Schwartz seminorm bound
for a tempered distribution then implies
$\mu([0,T])\leq C(1+T)^M$ for some $C,M$ and all $T\geq1$.
For $0\leq\sigma<1$, however,
\[
 \mu_\sigma([0,T])\geq e^{-\sigma T}\psi(e^T)
                  \geq c e^{(1-\sigma)T}
\]
for large $T$. For $\sigma<0$, use $\mu_\sigma\geq\mu_0$.
Both contradict polynomial mass growth.
\end{proof}

Multiplication by $e^{-\sigma t}$ is explicit here; it is not a
coordinate relabelling. Equivalently, for $\sigma\geq1$ the
unweighted $\mu_0$ acts continuously on
$\mathscr S_\sigma=e^{-\sigma t}\mathscr S(\mathbb R)$ with the
transported Schwartz topology. It acts on $C_c^\infty$ without any
choice of weight. The change $x=e^t$ pulls a compact smooth test
back by $f(t)=g(e^t)$ and gives the exact pairing
$\mu_\sigma(f)=\sum\Lambda(m)m^{-\sigma}g(m)$.
These are maps between specified test spaces, not Hilbert-space
intertwiners.

\begin{proposition}[Laplace and Fourier transforms on their domains]
\label{v4-arith:pm:laplace}
For $\Re s>1$,
\[
 \int e^{-st}\,d\mu_0(t)
 =\sum_{m\geq2}\frac{\Lambda(m)}{m^s}
 =-\frac{\zeta'(s)}{\zeta(s)}.
\]
With Fourier convention $\mathcal F_-f(u)=\int f(t)e^{-iut}dt$,
the finite measure $\mu_\sigma$, $\sigma>1$, has Fourier transform
$-\zeta'(\sigma+iu)/\zeta(\sigma+iu)$.
As $\sigma\downarrow1$ these transforms converge in the weak
tempered-distribution sense to $\mathcal F_-\mu_1$.
The unweighted and critically weighted measures do not have Fourier
transforms in $\mathscr S'$ by this construction.
\end{proposition}
\begin{proof}
The absolutely convergent Euler product for $\Re s>1$ may be
logarithmically differentiated on compact subsets there, giving
the displayed Dirichlet series. Absolute convergence also permits
integration of the finite weighted measure against the Fourier
kernel. For every Schwartz test, dominated convergence in the
seminorm estimate of Theorem~\ref{v4-arith:pm:weights} gives
$\mu_{1+\varepsilon}\to\mu_1$ weakly in $\mathscr S'$.
Fourier continuity proves the boundary assertion. Meromorphic
continuation of the right side does not change the abscissa of
absolute convergence of the original positive-measure integral.
\end{proof}

\subsection{The explicit formula on compact Fourier tests}
\label{v4-arith:pm:explicit-section}

For $f\in\mathscr A$, use the entire function
$h(z)=\int f(t)e^{izt}dt$, so that
$f(t)=(2\pi)^{-1}\int h(u)e^{-iut}du$.
This is the Fourier test carrier of
Definition~\ref{v4-arith:mc:test-algebra}.
For $\rho=\beta+i\gamma_\rho$ put
$z_\rho=-i(\rho-1/2)=\gamma_\rho-i(\beta-1/2)$.
Lemma~\ref{v4-arith:pm:zero-count} gives
$\#\{\rho:|\gamma_\rho|\leq T\}=O(T\log(T+2))$, with multiplicity.
It follows that $\mu_{\mathrm{ord}}=\sum_\rho\delta_{\gamma_\rho}$
is tempered. Since $0\leq\beta\leq1$, integration by parts gives rapid
decay of $h(z_\rho)$ uniformly in this horizontal strip.
Thus $\mathfrak Z_\xi(h)=\sum_\rho h(z_\rho)$ converges absolutely
and is continuous on $\mathscr A$ through $f\mapsto h$.
The estimates therefore follow from the theta integral without
requiring an asymptotic prime-counting formula.

\begin{theorem}[the full polar and archimedean terms]
\label{v4-arith:pm:weil}
For every $f\in C_c^\infty(\mathbb R)$ and its transform $h$ above,
\begin{align*}
 \mathfrak Z_\xi(h)
 &=h(i/2)+h(-i/2)
   +\frac1{2\pi}\int_{\mathbb R}
       h(u)\bigl(\Re\psi_\Gamma(1/4+iu/2)-\log\pi\bigr)\,du\\
 &\quad-\sum_{m\geq2}\frac{\Lambda(m)}{\sqrt m}
                  \bigl(f(\log m)+f(-\log m)\bigr),
\end{align*}
where $\psi_\Gamma=\Gamma'/\Gamma$ is the digamma function.
The prime sum is finite. In particular, it is evaluation of
$\mu_{1/2}$ on compact tests and does not assert temperedness of
that measure. The values $h(\pm i/2)$ are the polar terms, separate
from the real-line digamma density.
\end{theorem}
\begin{proof}
Set $H(s)=h(-i(s-1/2))=\int f(t)e^{(s-1/2)t}dt$ and choose $a>1$.
The residue theorem and $\xi(s)=\xi(1-s)$ give
\[
 \sum_\rho H(\rho)
 =\frac1{2\pi i}\int_{\Re s=a}
             (H(s)+H(1-s))\frac{\xi'(s)}{\xi(s)}\,ds.
\]
Here is a convergence justification. The zero count allows heights
$T_j\in[j,j+1]$ at distance at least $j^{-2}$ from every zero
ordinate: the excluded intervals have total length $O(\log j/j)$.
The paired logarithmic derivative in
Proposition~\ref{v4-arith:pm:constant} is polynomially bounded on
the horizontal segments at these heights. For $|\rho|\leq2T_j$
bound each denominator by the stated distance and sum
$O(T_j\log T_j)$ terms; for $|\rho|>2T_j$ use
$|s/\{\rho(s-\rho)\}|\leq C T_j/|\rho|^2$.
Repeated integration by parts makes $H$ decrease faster than any
power, uniformly on the fixed strip. The horizontal integrals
therefore vanish, and the zero sum converges absolutely.

On the right line decompose
\[
 \frac{\xi'}\xi(s)
 =\frac1s+\frac1{s-1}-\frac12\log\pi
       +\frac12\psi_\Gamma(s/2)+\frac{\zeta'}\zeta(s).
\]
Fourier inversion and the absolutely convergent Dirichlet series
give the negative prime sum from the last term. For the remaining
terms move the line to $\Re s=1/2$. The crossed pole at $s=1$
contributes $H(1)+H(0)$. On the new line the rational sum
$1/s+1/(s-1)$ is odd in the imaginary coordinate, while
$H(s)+H(1-s)$ is even; its integral is zero. Changing $u$ to $-u$
in half of the gamma integral gives precisely its real part in the
displayed formula. Digamma growth is $O(\log(|u|+2))$ on this
line: in
$\psi_\Gamma(z)=-\gamma+\sum_{n\geq0}((n+1)^{-1}-(n+z)^{-1})$,
split the sum at $n=|u|+1$ for $z=1/4+iu/2$. The initial sum
is $O(\log(|u|+2))$ and the tail is bounded by
$C(|u|+1)\sum_{n>|u|+1}n^{-2}=O(1)$.
Thus all resulting integrals converge absolutely.
\end{proof}

When $f$ is even, $h$ is even and the prime term is
$-2\mu_{1/2}(f)$. The functional
$h\mapsto h(i/2)+h(-i/2)$ cannot be silently included in a density
on the real line. Likewise, $\mathfrak Z_\xi(h)$ uses the full
complex numbers $z_\rho$, whereas
$\mu_{\mathrm{ord}}(h)=\sum h(\gamma_\rho)$ uses their real parts.
Under RH these are equal. No projection is used in the theorem.

\subsection{Centering and the critical exponent}
\label{v4-arith:pm:centered-section}

For real $\sigma$ define, initially on compact smooth tests,
\[
 \nu_\sigma=\mu_\sigma
       -\mathbf1_{[0,\infty)}(t)e^{(1-\sigma)t}\,dt.
\]
This subtracts an explicit density and is a different distribution.
For $\sigma\geq1$ it is tempered by
Theorem~\ref{v4-arith:pm:weights}. In particular its transform at
$\sigma=1$ is the weak boundary value
\[
 \mathcal F_-\nu_1
 =\lim_{\varepsilon\downarrow0}
    \left(-\frac{\zeta'(1+\varepsilon+iu)}
                 {\zeta(1+\varepsilon+iu)}
                    -\frac1{\varepsilon+iu}\right).
\]
The removed half-line density has Fourier transform
$\pi\delta_0-i\,\mathrm{pv}(1/u)$, by taking the real and imaginary
parts of $(\varepsilon+iu)^{-1}$. Thus centering specifies both the
Dirac and principal-value boundary terms.

\begin{theorem}[critical centering as a distributional criterion]
\label{v4-arith:pm:centered}
The distribution $\nu_{1/2}$ is tempered if and only if every
nontrivial zero of $\zeta$ has real part $1/2$.
This is a criterion for a specified centered distribution; it
constructs no arithmetic Hilbert pairing or occupation intertwiner.
\end{theorem}
\begin{proof}
First restrict a test $f$ to $C_c^\infty(0,\infty)$.
The digamma integral representation
\[
 \psi_\Gamma(z)=-\gamma+
  \int_0^\infty\frac{e^{-v}-e^{-zv}}{1-e^{-v}}\,dv
 \qquad(\Re z>0)
\]
is \href{https://dlmf.nist.gov/5.9.E16}{DLMF 5.9.16}.
This identity
and Fourier inversion, with $v=2t$, turn the archimedean integral
of Theorem~\ref{v4-arith:pm:weil} into
$-\int_0^\infty f(t)e^{-t/2}/(1-e^{-2t})\,dt$.
The constant terms vanish since $f(0)=0$; near $v=0$ the
subtracted representation is integrable, and compact support away
from zero justifies the Fourier inversion. The term $h(i/2)$
cancels the first exponential in this denominator. Hence, on
$(0,\infty)$,
\[
 \nu_{1/2}(f)
 =-\sum_\rho\int f(t)e^{(\rho-1/2)t}\,dt
   -\int_0^\infty\frac{e^{-t/2}}{e^{2t}-1}f(t)\,dt.
\]
Under RH the zero sum extends to a tempered distribution on the
whole real line by the exact map
$f\mapsto\mu_{\mathrm{ord}}(\mathcal F_+f)$, where
$\mathcal F_+f(u)=\int f(t)e^{iut}dt$.
Multiply this distribution by a smooth cutoff that vanishes near
zero and is one for $t\geq1$. The last displayed density then
decays exponentially. The difference from $\nu_{1/2}$ has compact
support and is a distribution, so it too is tempered. This proves
the forward construction under RH.

Conversely suppose $\nu_{1/2}\in\mathscr S'$. Its support is
contained in $[0,\infty)$. Choose a smooth function $\eta$ that
is one there and zero on $(-\infty,-1]$. For $\Re s>0$ define
$L(s)=\nu_{1/2}(\eta(t)e^{-st})$.
The tests and all their $s$-derivatives lie in $\mathscr S$, locally
uniformly for $\Re s>0$; hence $L$ is holomorphic there.
For $\Re s>1/2$ the convergent integral gives
\[
 L(s)=-\frac{\zeta'(s+1/2)}{\zeta(s+1/2)}
                         -\frac1{s-1/2}.
\]
Meromorphic uniqueness extends this equality to $\Re s>0$.
The pole at $s=1/2$ cancels, but any zero with $\Re\rho>1/2$
would leave a pole of nonzero integer residue at $s=\rho-1/2$.
Holomorphy excludes these zeros. The functional equation excludes
zeros with $\Re\rho<1/2$ as well.
\end{proof}

The unweighted centered distribution $\nu_0$ is not tempered.
Indeed the same Laplace argument would make
$-\zeta'(s)/\zeta(s)-1/(s-1)$ holomorphic for $\Re s>0$.
The infinitely many xi zeros of Lemma~\ref{v4-arith:pm:zero-count},
together with reflection in $\Re s=1/2$, provide a zero with
positive real part, and hence an uncancelled pole.

\begin{corollary}[the ordinate shadow on the full test class]
\label{v4-arith:pm:shadow}
Equality $\mathfrak Z_\xi(h)=\mu_{\mathrm{ord}}(h)$ for every
$h(z)=\int f(t)e^{izt}dt$, $f\in C_c^\infty(\mathbb R)$,
is equivalent to RH.
\end{corollary}
\begin{proof}
Under RH the equality is termwise. Conversely the equality makes
the zero functional on compact Fourier tests the restriction of a
tempered Fourier transform. The cutoff argument in
Theorem~\ref{v4-arith:pm:centered} then proves
$\nu_{1/2}\in\mathscr S'$. That theorem implies RH.
\end{proof}

The uncentered critical measure remains non-tempered even under RH.
Centering subtracts its exponentially growing main density; a
statement about this difference does not change the carrier of the
original prime measure. These distinctions are required before
any Fourier, Mellin, or spectral comparison can be made.
